\chapter{Documentation des sous-modèles}
\label{ch:sousmodeles}

Ce chapitre documente chaque mécanisme du modèle selon une structure fixe : \textbf{A} ce que le sous-modèle
représente ; \textbf{B} les variables ; \textbf{C} les hypothèses ; \textbf{D} les équations, telles qu'elles sont
codées ; \textbf{E} la calibration ; \textbf{F} les limites. Les équations reprennent les notations du code
(fichiers \code{src/m7\_equations.py} et \code{src/m8\_modele.py}). Les coefficients sont ceux de la version 0.6 ;
les écarts-types figurent entre parenthèses. Conventions : $\ln$ est le logarithme népérien ; $\Delta \ln X_t =
\ln X_t - \ln X_{t-1}$ est, à peu de chose près, le taux de croissance de $X$ ; l'indice $b$ désigne une branche,
$t$ une année, $B = 2022$ l'année de base.

\begin{encadre}{Encadré 2.1 — Lire une équation « à correction d'erreur »}
La plupart des équations de comportement ont la même forme :
\[
\Delta \ln Y_t = \underbrace{a\,\Delta \ln X_t}_{\text{réaction immédiate}}
\;+\; \underbrace{\lambda\,\bigl(\ln Y_{t-1} - \ln Y^{*}_{t-1}\bigr)}_{\text{rappel vers l'équilibre}} \;+\; c ,
\qquad \lambda < 0 .
\]
Le premier terme dit comment $Y$ réagit \emph{la première année} à une variation de $X$ ; le second ramène
progressivement $Y$ vers sa valeur d'équilibre $Y^{*}$ (par exemple, la consommation vers une part stable du revenu).
Si $\lambda = -0{,}22$, l'écart à l'équilibre se réduit de 22~\% chaque année : il en reste la moitié au bout de
trois ans environ ($0{,}78^3 \approx 0{,}47$). Ce type d'équation, dit « à correction d'erreur » (MCE), est la forme
standard des modèles macro-économétriques depuis les travaux d'Engle et Granger (1987).
\end{encadre}

\section{Démographie et population active (M1, M9)}
\label{sec:m1}

\paragraph{A. Objet.} La démographie fixe deux grandeurs décisives pour l'économie : le nombre de personnes en âge
de travailler (et donc l'offre de travail) et le nombre de retraités (et donc la masse des pensions). Le rapport entre
les deux conditionne la soutenabilité du système social.

\paragraph{B. Variables.} $P_{s,a,t}$ : population au 1\textsuperscript{er} janvier par sexe $s$ et âge révolu
$a$ (0 à 120 ans) ; $q_{s,a,t}$ : quotients de mortalité ; $f_{a,t}$ : taux de fécondité ; $SM_{s,a,t}$ : solde
migratoire ; $\tau_{s,a}$ : taux d'activité ; $AOD_g$ : âge d'ouverture des droits à la retraite de la génération $g$.
Source : Insee, projections de population 2026 \src{Insee, projections 2026} ; Eurostat \code{lfsa\_argan} (taux d'activité).

\paragraph{C. Hypothèses.} Le module applique les hypothèses centrales de l'Insee : fécondité de 1,45 enfant par
femme à partir de 2028, solde migratoire de +150\,000 personnes par an, gains d'espérance de vie au rythme 2005-2025
(89,5 ans pour les femmes et 86,7 ans pour les hommes en 2070). Les taux d'activité par sexe et âge avant 55 ans
restent à leur niveau de 2024 ; ceux de 55 à 74 ans dépendent de l'âge légal (M9).

\paragraph{D. Équations.} Méthode des composantes (diagramme de Lexis) :
\begin{align}
P_{s,A,t+1} &= P_{s,A-1,t}\,(1-q_{s,A,t}) + SM_{s,A,t}\,\bigl(1-\tfrac{q_{s,A,t}}{2}\bigr) + AJ_{s,A,t} \\
N_t &= \sum_{A=15}^{50} f_{A,t}\,\frac{P^F_{A-1,t}+P^F_{A,t+1}}{2},\qquad P_{s,0,t+1} = N_t\,\pi_s\,(1-q_{s,0,t})+\dots \\
\text{Actifs}_t &= \sum_{s}\sum_{a} P_{s,a,t}\,\tau_{s,a}
\end{align}
avec $\pi_F = 100/205$ la part des filles à la naissance et $AJ$ l'ajustement statistique de l'Insee (nul après 2028).
L'identité $P_{t+1}-P_t = N_t - D_t + SM_t + AJ_t$ est contrôlée à chaque pas (tolérance 0,001 personne).

Pour les seniors (M9), le taux d'activité de 55 à 74 ans est exprimé relativement à celui des 50-54 ans et dépend de
l'écart entre l'âge et l'$AOD$ de la génération. Avant l'$AOD$ : sorties précoces en rampe sur 5 ans (carrières
longues, invalidité, chômage) ; après l'$AOD$ : une part reste active puis décroît exponentiellement. Un paramètre
$\beta$ mesure la part du profil qui se décale avec l'âge légal (la part $1-\beta$ reste attachée à 62 ans).

\begin{exemple}
Relever l'âge légal d'un an décale environ les trois quarts ($\beta \approx 0{,}75$) des départs d'un an. Sur une
génération de 800\,000 personnes dont 60~\% sont encore actives à 62 ans, cela maintient environ
$0{,}75 \times 0{,}6 \times 800\,000 \approx 360\,000$ personnes-années d'activité supplémentaires par génération concernée.
Avec $\beta = 1$ (décalage complet), l'effet serait d'un tiers plus fort ; avec $\beta = 0{,}5$, d'un tiers plus faible.
\end{exemple}

\paragraph{E. Calibration.} Démographie : aucun paramètre calé, les hypothèses publiées sont appliquées ; le module
reproduit les 29 scénarios de l'Insee à moins de $3{,}4\cdot10^{-5}$ près sur la population totale 2024-2070
\preuve{A}. Activité des seniors : estimation par moindres carrés non linéaires sur l'enquête Emploi 2003-2019 ;
$\beta = 0{,}78$ (femmes) et $0{,}75$ (hommes), erreur quadratique moyenne 1,0 à 1,1 point \preuve{E}. Hors
échantillon (2020-2025, relèvement de l'AOD de 62 ans à 62 ans et 9 mois), l'écart est de 0 à 3 points pour les
60-64 ans. Sur 2010-2019, neuf seniors maintenus actifs sur dix étaient en emploi (rapport $\Delta$ taux
d'emploi / $\Delta$ taux d'activité = 0,89).

\paragraph{F. Limites.} Pas de durée d'assurance, de décote ni de surcote ; taux d'activité avant 55 ans figés ; la
fécondité et les migrations ne réagissent pas à l'économie ; l'enquête Emploi a connu une rupture de série en 2021.

\section{Production, branches et échanges entre produits (M5)}
\label{sec:m5}

\paragraph{A. Objet.} Produire une voiture demande de l'acier, de l'électricité, du transport, des services : chaque
branche achète aux autres. Le tableau entrées-sorties décrit ces échanges et permet de calculer combien de production,
d'importations et d'emplois sont mobilisés, directement et indirectement, par un euro de demande finale.

\paragraph{B. Variables.} $x$ : production par produit (64 produits, nomenclature A64) ; $f^d$ : demande finale
adressée à la production intérieure ; $A^d$, $A^m$ : coefficients techniques des consommations intermédiaires
intérieures et importées ; $a^v$ : valeur ajoutée par unité produite ; $p$ : prix de production ; $p^m$ : prix
d'importation ; $\varphi$ : facteur commun d'ajustement des parts importées. Source : Eurostat
\code{naio\_10\_cp1700} (tableau symétrique, 2010-2022), \code{naio\_10\_cp15} (prix d'importation).

\paragraph{C. Hypothèses.} Coefficients techniques fixes à l'année de base 2022 : pas de substitution entre
intrants ni entre produit français et importé, sauf par le facteur commun $\varphi$ qui ajuste toutes les parts
importées (sauf l'énergie, non concurrente) pour respecter l'équation agrégée d'importations (section~\ref{sec:commerce}).

\paragraph{D. Équations.}
\begin{align}
x &= \bigl(I - A^d(\varphi)\bigr)^{-1} f^d(\varphi), \qquad m = A^m(\varphi)\,x + f^m(\varphi) \\
p &= \bigl(p^m A^m + p^v a^v\bigr)\bigl(I - A^d\bigr)^{-1}
\end{align}
où $p^v$ est le prix de la valeur ajoutée unitaire de la branche productrice (salaires, marges, impôts sur la
production). La valeur ajoutée de la branche $b$ en volume est $Q_b = \sum_{p\in b} a^v_p x_p$.

\begin{exemple}
Le contenu en importations de la demande finale intérieure est de 0,18~€ par euro en 2021 ; il atteint 0,73 pour les
produits raffinés, 0,52 pour l'aéronautique et les autres matériels de transport, 0,50 pour la métallurgie et 0,48 pour
l'automobile \src{validation M5}. Une hausse de 50~\% des prix de l'énergie importée, à marges fixes, élève les prix de
production de 1,3~\% et les prix à la consommation de 1,7~\%.
\end{exemple}

\paragraph{E. Calibration.} Observé \preuve{A/R} : identités vérifiées à 2\,M€ près pour 2010-2022. Test des prix
2012-2019 avec les prix de valeur ajoutée observés : erreur de 0,3 à 0,8 point par produit, corrélation de 0,89 à 0,98.

\paragraph{F. Limites.} En 2022, le modèle de prix à coefficients fixes surestime fortement la hausse de certains
prix (raffinage +74~\% prévu contre +43~\% observé), faute de substitutions et de variations de marges ; le tableau
2022 est le dernier publié (2023-2024 vides chez Eurostat).

\section{Emploi et chômage}
\label{sec:emploi}

\paragraph{A. Objet.} Les entreprises ajustent leurs effectifs à leur production, avec retard, et selon le coût réel
du travail. Le chômage est la part de la population active sans emploi.

\paragraph{B. Variables.} $N_{b,t}$ : emploi de la branche (personnes, milliers) ; $Q_{b,t}$ : valeur ajoutée en
volume ; $CN_{b,t}$ : coût salarial par tête ; $PVA_{b,t}$ : prix de la valeur ajoutée ; $g_b$ : tendance de
productivité de la branche ; $u_t$ : taux de chômage (France métropolitaine, BIT) ; $PA_t$ : population active.
Sources : Eurostat \code{nama\_10\_a64}, \code{nama\_10\_a64\_e} ; Insee, série 001688526.

\paragraph{C. Hypothèses.} Rendements d'échelle constants à long terme (l'emploi est proportionnel à la production,
corrigé de la productivité) ; élasticité de substitution entre capital et travail $\sigma = 0{,}4$ ; les emplois non
marchands sont proportionnels au volume de la production non marchande ; la productivité tendancielle par branche est
celle des huit dernières années observées.

\paragraph{D. Équations.} Pour chaque branche marchande :
\begin{multline}
\Delta\ln N_{b,t} = c_b + \underset{(0{,}043)}{0{,}259}\,\Delta\ln Q_{b,t} + \underset{(0{,}018)}{0{,}012}\,\Delta\ln Q_{b,t-1}
+ \underset{(0{,}057)}{0{,}317}\,\Delta\ln N_{b,t-1} \\
\underset{(0{,}031)}{-\,0{,}206}\Bigl[\ln N_{b,t-1} - \ln Q_{b,t-1} + \sigma\,\ln\frac{CN_{b,t-1}}{PVA_{b,t-1}} + g_b\,(t-1-B)\Bigr]
\end{multline}
\begin{equation}
u_t = 100\left(1 - \frac{\sum_b N_{b,t} + N^{nm}_t}{PA_t}\right)
\end{equation}
Le crochet est l'écart entre l'emploi et son niveau d'équilibre $\ln N^* = \ln Q - \sigma \ln(CN/PVA) - g_b\,t$.

\begin{exemple}
Si la production d'une branche augmente durablement de 1~\%, l'emploi augmente de 0,26~\% la première année, puis de
0,26~\% $\times$ 0,317 par inertie et de 21~\% de l'écart restant par le rappel : environ +0,5~\% au bout de deux ans,
+0,9~\% au bout de six. Si le coût réel du travail augmente de 1~\%, l'emploi d'équilibre baisse de $\sigma = 0{,}4$~\%.
Avec $\sigma = 0{,}2$ (bas de la littérature), une hausse des salaires détruit deux fois moins d'emplois ; avec
$\sigma = 0{,}7$, presque deux fois plus.
\end{exemple}

\paragraph{E. Calibration.} Estimé sur 25 branches, 1997-2019, effets fixes et tendances de branche \preuve{E}.
L'élasticité de long terme à la production est imposée à 1 (la version libre donne 0,44, biaisée par la colinéarité
avec les tendances). $\sigma = 0{,}4$ est \preuve{C} : non identifié sur les données (estimation libre $-$0,21, non
robuste), fixé au milieu des estimations françaises (0,2 à 0,7).

\paragraph{F. Limites.} Une seule élasticité pour tous les niveaux de qualification : le modèle ne représente pas la
sensibilité plus forte de l'emploi peu qualifié au coût du travail, documentée par les évaluations des allègements de
cotisations (Kramarz et Philippon, 2001). Il sous-estime donc probablement l'effet des allègements ciblés et l'effet
d'une hausse du SMIC sur l'emploi peu qualifié. Pas de flux d'embauches et de licenciements, pas d'appariement, pas
d'automatisation au-delà de la substitution capital-travail ($\sigma$).

\section{Formation des salaires}
\label{sec:salaires}

\paragraph{A. Objet.} Le salaire moyen s'ajuste aux prix (indexation, négociations), à la productivité, et à l'état du
marché du travail (un chômage élevé modère les salaires : courbe de Phillips).

\paragraph{B. Variables.} $W_t$ : salaire moyen par tête du secteur marchand ; $PC_t$ : prix de la consommation ;
$PROD_t$ : productivité par tête ; $PVA^m_t$ : prix de la valeur ajoutée marchande ; $u^*_t$ : chômage structurel
(7,5~\% en projection).

\paragraph{D. Équation.}
\begin{multline}
\Delta\ln W_t = c + \underset{(0{,}222)}{0{,}394}\,\Delta\ln PC_t + \underset{(0{,}120)}{0{,}032}\,\Delta\ln PROD_t \\
\underset{(0{,}057)}{-\,0{,}093}\bigl[\ln W_{t-1} - \ln PVA^m_{t-1} - \ln PROD_{t-1}\bigr] \;-\; 0{,}004\,(u_{t-1}-u^*_t)
\end{multline}
Le crochet est (le logarithme de) la part des salaires dans la valeur ajoutée : si elle est au-dessus de sa norme, les
salaires ralentissent.

\begin{exemple}
Une inflation de 2~\% augmente les salaires de 0,8~\% la même année ; le reste est rattrapé à raison de 9~\% de l'écart
par an, si bien que la part des salaires revient à sa norme en une dizaine d'années. Un point de chômage au-dessus du
niveau structurel freine la croissance des salaires de 0,4 point par an.
\end{exemple}

\paragraph{E. Calibration.} Indexation et rappel estimés \preuve{E} (l'indexation a un $t$ de 1,8 seulement).
Effet du chômage \preuve{C} : l'estimation libre donne un effet positif (+0,0066), biaisé par la composition de
l'emploi (les emplois à bas salaire disparaissent d'abord en récession, ce qui augmente mécaniquement le salaire
moyen) ; la valeur retenue ($-$0,4~\% par point) est dans la fourchette des estimations françaises.

\paragraph{F. Limites.} Un seul salaire moyen : pas de distinction entre SMIC et autres salaires, ni entre secteurs
exposés et abrités. Une hausse du SMIC est représentée par sa diffusion supposée à la masse salariale (0,11 pour 1,
hypothèse \preuve{H}) et par son coût public direct (allègements indexés sur le SMIC : 0,66\md{} par point).

\section{Prix et inflation}
\label{sec:prix}

\paragraph{A. Objet.} Les entreprises fixent leurs prix en appliquant une marge à leurs coûts ; les hausses de prix
importés (énergie notamment) se transmettent avec retard.

\paragraph{D. Équations.} Prix de la valeur ajoutée de chaque branche, en fonction du coût salarial unitaire
$CSU_b = CN_b N_b / Q_b$ :
\begin{multline}
\Delta\ln PVA_{b,t} = c_b + \underset{(0{,}113)}{0{,}670}\,\Delta\ln CSU_{b,t} + \underset{(0{,}072)}{0{,}109}\,\Delta\ln Q_{b,t}
+ \underset{(0{,}001)}{0{,}0015}\,TUCD_{b,t} \\ \underset{(0{,}057)}{-\,0{,}117}\bigl[\ln PVA_{b,t-1} - \ln CSU_{b,t-1} - mk_{b,t-1}\bigr] + \Delta mk_{b,t}
\end{multline}
où $TUCD$ est l'écart du taux d'utilisation des capacités à sa moyenne (branches industrielles) et $mk$ un choc de
marge éventuel (impôts de production, droits de douane). Transmission des prix d'importation, par produit :
\begin{equation}
\ln p^{m,\text{eff}}_t = \ln p^{m,\text{eff}}_{t-1} + \theta\,\bigl(\ln p^m_t - \ln p^{m,\text{eff}}_{t-1}\bigr),\qquad \theta = \underset{(\text{estimé})}{0{,}495}.
\end{equation}
Les prix de la consommation, de l'investissement et des exportations sont ensuite les moyennes, pondérées par les
structures du tableau entrées-sorties, des prix de production $p$ (section~\ref{sec:m5}) et des prix importés
répercutés.

\begin{exemple}
Une hausse de 10~\% des prix de l'énergie importée n'est répercutée qu'à moitié la première année ($\theta \approx 0{,}5$),
aux trois quarts la deuxième. Une hausse de 1~\% du coût salarial unitaire relève le prix de la valeur ajoutée de
0,67~\% la première année, puis entièrement à long terme (homogénéité imposée).
\end{exemple}

\paragraph{E. Calibration.} \preuve{E} sur 1997-2019 pour les prix de valeur ajoutée (la version libre donne une
transmission de long terme de 0,74 au lieu de 1 imposé) ; $\theta$ estimé sur 1998-2024 : sur 1998-2019, la
transmission était immédiate ; le délai n'apparaît qu'avec l'épisode énergétique de 2021-2024 (tarifs réglementés,
bouclier tarifaire, contrats).

\paragraph{F. Limites.} Pas de prix de marché de l'électricité ni de mécanisme de reconstitution de marges ; le
modèle a sous-estimé les prix de 2,8~\% sur 2023-2024 (validation M8).

\section{Consommation et épargne des ménages (M7, M14)}
\label{sec:conso}

\paragraph{A. Objet.} La consommation des ménages (la moitié du PIB) dépend avant tout de leur revenu réel. Les
ménages modestes dépensent une part plus grande d'un euro supplémentaire que les ménages aisés.

\paragraph{B. Variables.} $C_t$ : consommation en volume ; $RDB_t$ : revenu disponible brut (M€ courants) ;
$RDB^{\text{eff}}_t$ : revenu « effectif » corrigé de l'hétérogénéité ; $PC_t$ : prix de la consommation ; taux
d'épargne $s_t = (RDB_t - PC_t C_t)/RDB_t$.

\paragraph{D. Équations.}
\begin{equation}
\Delta\ln C_t = c + \underset{(0{,}205)}{0{,}853}\,\Delta\ln\frac{RDB^{\text{eff}}_t}{PC_t} \;\underset{(0{,}211)}{-\,0{,}218}\Bigl[\ln C_{t-1} - \ln\frac{RDB^{\text{eff}}_{t-1}}{PC_{t-1}}\Bigr]
\end{equation}
\begin{equation}
RDB^{\text{eff}}_t = RDB_t + \sum_{d=1}^{10} \Delta R_{d,t}\Bigl(\frac{\pi_d}{\bar\pi}-1\Bigr) + \Delta Div^{\text{mén}}_t\Bigl(\frac{\pi_{\text{top 1}}}{\bar\pi}-1\Bigr)
\end{equation}
où $\Delta R_d$ est la variation du revenu du décile $d$ due à la mesure (microsimulation), $\pi_d$ la propension
marginale à consommer du décile (0,90 ; 0,85 ; 0,80 ; 0,75 ; 0,72 ; 0,68 ; 0,64 ; 0,60 ; 0,53 ; 0,40), $\bar\pi$ sa moyenne
pondérée par le revenu, et $\pi_{\text{top 1}} = 0{,}25$ la propension du dernier centile, appliquée aux dividendes.

\begin{exemple}
Un transfert de 1\md{} aux ménages du premier décile compte pour $1 \times 0{,}90/\bar\pi \approx 1{,}38$\md{} dans
l'équation de consommation ($\bar\pi \approx 0{,}65$) ; une baisse d'impôt de 1\md{} pour le dixième décile ne compte
que pour $0{,}40/0{,}65 \approx 0{,}62$\md. À montant égal, la première mesure soutient donc environ deux fois plus la
consommation la première année. Seul le \emph{profil relatif} des propensions intervient : la propension moyenne reste
celle de l'équation estimée.
\end{exemple}

\paragraph{E. Calibration.} Équation agrégée \preuve{E} (1997-2019) ; l'élasticité de court terme varie de 0,54 à 1,15
selon la période d'estimation et la force de rappel n'est pas significative ($t = -1{,}0$). Profil des propensions
\preuve{C} : cohérent avec les taux d'épargne par quintile de l'Insee (de $-$29~\% pour le premier quintile à 27~\% pour le
cinquième en 2022) \src{Insee Focus n\textsuperscript{o} 338, 2024} et avec Garbinti et Lamarche (2014) ; variantes
« plat » et « accentué » en sensibilité.

\paragraph{F. Limites.} Pas d'effet de richesse, de taux d'intérêt, d'épargne de précaution ni d'anticipation
(pas d'« équivalence ricardienne ») ; pas de distinction entre revenu permanent et temporaire (une prime ponctuelle a
probablement un effet d'un tiers plus faible la première année que celui simulé) ; le taux d'épargne de référence
(17,9~\% en 2025, ramené à 15,5~\% à long terme) est une hypothèse de projection.

\section{Investissement et accumulation du capital}
\label{sec:invest}

\paragraph{A. Objet.} Les entreprises investissent pour répondre à la demande (effet d'accélérateur) et d'autant plus que
le coût du capital (taux d'intérêt, impôt sur les sociétés) est bas. Les ménages investissent dans le logement selon
leur revenu. Le capital installé détermine les capacités de production.

\paragraph{D. Équations.} Investissement de la branche $b$ (hors effet du coût du capital), puis effet du coût d'usage
$CU$ appliqué en niveau :
\begin{multline}
\Delta\ln \tilde I_{b,t} = c_b + \underset{(0{,}204)}{0{,}670}\,\Delta\ln Q_{b,t} + \underset{(0{,}097)}{0{,}256}\,\Delta\ln Q_{b,t-1}
\underset{(0{,}109)}{-\,0{,}285}\ln \tilde I_{b,t-1} + \underset{(0{,}096)}{0{,}192}\ln K_{b,t-1} + \underset{(0{,}048)}{0{,}113}\ln Q_{b,t-1}
\end{multline}
\begin{equation}
CU_t = \frac{r_t + \delta}{1-\tau^{IS}_t},\qquad \ln m_t = \ln m_{t-1} + 0{,}285\,\bigl(\varepsilon\,[\ln CU_t - \ln CU_B] - \ln m_{t-1}\bigr),\qquad I_{b,t} = \tilde I_{b,t}\, m_t
\end{equation}
avec $\delta = 0{,}07$, $\varepsilon = -0{,}184$ (élasticité de long terme estimée sur l'investissement des sociétés
non financières, écart-type de $\lambda$ 0,14). Capital : $K_{b,t} = (1-\delta_b)K_{b,t-1} + I_{b,t} + I^{\text{cap}}_{b,t}$.
Logement des ménages :
\begin{equation}
\Delta\ln I^{\text{mén}}_t = c + \underset{(0{,}35)}{1{,}48}\,\Delta\ln\frac{RDB_t}{PC_t} \underset{(0{,}20)}{-\,0{,}282}\Bigl[\ln I^{\text{mén}}_{t-1}-\ln\frac{RDB_{t-1}}{PC_{t-1}}\Bigr]
\end{equation}
L'investissement public est exogène : chaque milliard supplémentaire s'ajoute à la demande et, par une élasticité de
la production au capital public de 0,08 (capital public net égal à 55~\% du PIB, dépréciation de 4~\% par an,
productif l'année suivant la dépense), à la productivité.

\begin{exemple}[Exemple : baisse de l'impôt sur les sociétés de 25 \% à 20 \%]
$1-\tau$ passe de 0,75 à 0,80 : le coût d'usage baisse de $\ln(0{,}80/0{,}75) = 6{,}5~\%$. L'investissement de long
terme augmente de $0{,}184 \times 6{,}5 \approx 1{,}2~\%$, dont 28~\% la première année (0,34~\%). Avec
$\varepsilon = -0{,}4$ (haut de la littérature internationale), l'effet doublerait ; avec $\varepsilon = -0{,}1$, il
serait presque divisé par deux. La recette d'IS perdue (environ 12\md{} de gain statique pour les entreprises) n'est
pas récupérée par l'activité supplémentaire (chapitre~\ref{ch:experiences}).
\end{exemple}

\paragraph{E. Calibration.} Accélérateur \preuve{E} ; élasticité au coût d'usage \preuve{E, fragile} (23 observations,
prévision hors échantillon médiocre) ; élasticité au capital public \preuve{L} (méta-analyse de Bom et Ligthart, 2014 ;
plage 0,04-0,12).

\paragraph{F. Limites.} L'effet du taux de marge sur l'investissement, estimé positif mais non significatif ($t=1{,}4$),
n'est pas utilisé ; pas de contrainte de financement ni d'effet de trésorerie ; aucune dépense publique autre que
l'investissement (éducation, recherche, santé) n'a d'effet sur la productivité ; l'investissement de capacité de la
réindustrialisation est décidé de façon exogène.

\section{Productivité, capacités et croissance potentielle}
\label{sec:fermeture}

\paragraph{A. Objet.} À long terme, la production est limitée par la quantité de travail et de capital disponibles et par
leur efficacité. Un modèle tiré par la demande doit donc être « fermé » par un mécanisme qui ramène l'économie vers son
potentiel : politique monétaire, crédit, compétitivité.

\paragraph{C. Hypothèses.} (1) La productivité tendancielle de chaque branche est un paramètre de scénario (moyenne
des 8 dernières années observées ; 0,7~\% par tête pour l'ensemble de l'économie après 2031, hypothèse centrale du COR).
(2) Le taux d'utilisation des capacités de l'industrie influence (faiblement) les prix. (3) \textbf{Fermeture d'offre}
\preuve{H} : un supplément de demande intérieure $Z$ (60~\% consommation, 40~\% investissement) ramène progressivement
le chômage vers son niveau structurel $u^*$.

\paragraph{D. Équations.}
\begin{equation}
Z_t = Z_{t-1} + \Bigl[\rho\,(u_{t-1}-u^*) + 0{,}3\,\bigl((u_{t-1}-u^*)-(u_{t-2}-u^*)\bigr)\Bigr]\times\frac{1{,}5}{100},\qquad \rho = 0{,}15
\end{equation}
$Z$ est exprimé en part de la production totale ; 1,5 est le coefficient d'Okun (points de PIB par point de chômage).
C'est un correcteur « proportionnel-intégral » : tant que le chômage est au-dessus de $u^*$, la demande supplémentaire
s'accumule. Taux d'utilisation des capacités : $TUC_{b,t} = TUC_{b,B}\,(Q_{b,t}/Q_{b,B})/(K_{b,t-1}e^{g^k_b(t-B)}/K_{b,B-1})$.

\begin{exemple}
Une mesure qui augmente le chômage de 0,3 point déclenche chaque année un supplément de demande de
$0{,}15\times0{,}3\times1{,}5 = 0{,}07$ point de PIB, qui s'accumule : au bout de cinq ans, environ 0,25 point de PIB, ce
qui ramène le chômage vers sa valeur de départ. Avec $\rho = 0$ (vision « keynésienne » : rien ne ramène l'économie
au plein emploi), les effets de demande sont permanents ; avec $\rho = 0{,}4$, ils disparaissent en trois à quatre ans.
\end{exemple}

\paragraph{E. Calibration.} $\rho$ n'est pas estimé : il est choisi pour que les multiplicateurs de dépense publique
décroissent comme dans Mésange et FR-BDF (1,2 la première année, 0,85 la deuxième, 0,6 à cinq ans). \textbf{C'est le
premier facteur d'incertitude du modèle} sur le PIB et l'emploi à moyen terme (criblage de Morris, contre-expertise
v0.6).

\paragraph{F. Limites.} Mécanisme réduit et symétrique ; pas de politique monétaire nationale (zone euro) autre que
celle du module M11 ; pas d'hystérèse (un chômage durablement élevé ne dégrade pas le potentiel) ; la productivité ne
dépend ni de l'éducation, ni de la recherche, ni de l'automatisation, ni du temps de travail au-delà d'une hypothèse de
compensation (section~\ref{sec:catalogue}).

\section{Commerce extérieur, énergie et dépendances}
\label{sec:commerce}

\paragraph{A. Objet.} Les importations augmentent avec la demande (surtout l'investissement et les exportations, riches
en contenu importé) ; les exportations dépendent de la demande mondiale et de la compétitivité-prix.

\paragraph{D. Équations.} Importations totales en volume :
\begin{multline}
\Delta\ln M_t = c + \underset{(0{,}092)}{1{,}465}\,\Delta\ln D_t \underset{(0{,}101)}{-\,0{,}114}\,\Delta\ln\frac{PM^{nc}_t}{PPIB_t} \\
\underset{(0{,}149)}{-\,0{,}661}\Bigl[\ln M_{t-1} - 1{,}370\ln D_{t-1} + 0{,}163\ln\frac{PM^{nc}_{t-1}}{PPIB_{t-1}}\Bigr] + 0{,}0039\,t
\end{multline}
avec $D_t = 0{,}13\,C_t + 0{,}06\,G_t + 0{,}23\,(I_t+S_t) + 0{,}33\,X_t$ (demande pondérée par le contenu en importations,
en indices) et $PM^{nc}$ le prix des importations hors énergie (non concurrente). Exportations :
\begin{equation}
\Delta\ln X_t = c + \underset{(0{,}094)}{0{,}711}\,\Delta\ln DM_t \underset{(0{,}171)}{-\,0{,}595}\,\Bigl(\Delta\ln PX_t - \Delta\ln PW_t\Bigr)
\end{equation}
où $DM$ est la demande mondiale adressée (importations de l'Union européenne) et $PW$ le prix des concurrents. Le total
importé est réparti entre produits par le facteur $\varphi$ (section~\ref{sec:m5}) ; l'énergie (produits B, C19, D)
est non concurrente : sa part importée ne varie pas.

\begin{exemple}
Une hausse de 1~\% de la demande pondérée augmente les importations de 1,47~\% la première année. Une perte de
compétitivité-prix de 1~\% réduit les exportations de 0,6~\%, mais les importations de seulement 0,11~\% : le modèle
représente mieux la perte de marchés à l'exportation que la conquête du marché intérieur par les importations.
\end{exemple}

\paragraph{E. Calibration.} \preuve{E}, 1997-2019. L'élasticité-prix des importations ($-$0,11 à court terme, $-$0,16 à
long terme) est non significative et plus faible que la littérature ($-$0,3 à $-$0,6). Celle des exportations est
conforme aux estimations françaises ($-$0,4 à $-$0,8).

\paragraph{F. Limites.} Pas de taux de change hors zone euro ; pas de traitement spécifique des grands contrats
(aéronautique), ce qui explique l'erreur de 2023 ; dépendances matérielles représentées seulement en valeur (pas de
métaux critiques, pas de goulets physiques) ; pas de substitution entre énergies ni d'effet de la taxe carbone sur les
émissions.

\section{Revenus des ménages, profits et revenus versés au reste du monde}
\label{sec:revenus}

\paragraph{D. Équations.} Le revenu disponible brut est la somme de composantes indexées :
\begin{equation}
RDB_t = \underbrace{D1_B\,\frac{W_t}{W_B}\frac{N_t}{N_B}}_{\text{salaires}} + \underbrace{EBE^{mén}_t + D4_t}_{\text{patrimoine, indépendants}}
+ \underbrace{D62_t}_{\text{prestations}} + \text{autres}_t - \underbrace{D61_t}_{\text{cotisations}} - \underbrace{D5_t}_{\text{impôts}}
\end{equation}
Les prestations sont indexées sur les prix de l'année précédente et sur un volume exogène (pensions par M9 ; autres
prestations +1,4~\% par an, rythme 2011-2019). Les cotisations sont proportionnelles aux salaires, les impôts à la base
imposable (salaires, patrimoine, prestations). Les revenus du patrimoine et l'excédent des ménages suivent des équations
estimées sur le PIB nominal, les prix et le taux long :
\begin{equation}
\Delta\ln D4_t = c + 0{,}874\,\Delta\ln PIBN_t + 3{,}05\,\Delta\ln PC_t - 0{,}366\,[\ln D4_{t-1}-\ln PIBN_{t-1}] + 0{,}021\,r_{t-1}
\end{equation}
Profits (M14) : un supplément de profit après impôt $\Pi$ (baisse d'IS ou d'impôts de production) est distribué à
50~\% (plage 30-70~\%) ; 30~\% des dividendes vont aux non-résidents (plage 20-50~\%), 35~\% du reste directement aux
ménages résidents ; le solde est conservé par les entreprises. Le revenu national brut est $RNB = PIBN + R^{prim}$,
où $R^{prim}$ (revenus primaires nets reçus du reste du monde) vaut 3,1~\% du PIB à l'année de base, diminué des
dividendes supplémentaires versés à l'étranger.

\paragraph{F. Limites.} En 2023-2024, le RDB est sous-estimé de 7,9~\% en rétrosimulation : les revenus de la propriété
(+15~\% par an avec la hausse des taux) ne suivent pas le PIB nominal ; le niveau de solde extérieur du modèle
n'est pas fiable (chapitre~\ref{ch:reference}).

\section{Fiscalité des ménages (M4, M10, M13)}
\label{sec:fiscmen}

\paragraph{A. Objet.} Calculer, ménage par ménage, l'impôt sur le revenu, la CSG et les prestations, puis en déduire
la recette publique, la variation de revenu par décile et l'effet sur la pauvreté.

\paragraph{B. Données.} Échantillon de 2~\% des ménages ordinaires du fichier détail du recensement 2022 (163\,826
ménages, 344\,527 personnes, 201\,821 foyers fiscaux) ; revenus attribués par rang à partir de la distribution du revenu
fiscal de référence publiée par la DGFiP (IRCOM 2025, revenus 2024, 25 tranches) ; statistiques nationales par case de
déclaration (heures supplémentaires, revenus de l'étranger, épargne retraite) ; législation par OpenFisca-France.

\paragraph{C. Hypothèses.} Rang de revenu tiré de grilles publiées (salaires par catégorie socioprofessionnelle et sexe,
pensions) ; revenus du capital concentrés vers le haut ; non-recours aux prestations constant (taux apparents mesurés :
RSA 0,68 ; prime d'activité 0,66).

\paragraph{D. Calcul.} Pour une réforme $R$ : $\Delta T = \sum_h w_h\,[T_h(R) - T_h(\text{actuel})]$, ramené aux 41,6
millions de foyers de la DGFiP ; variation par décile de niveau de vie
$nv_h = (RFR_h - IR_h + \text{Prest}_h)/UC_h$ (échelle de l'OCDE modifiée). Les mesures dont l'assiette n'est pas
simulée ménage par ménage (TVA, taxe carbone, CSG sur les pensions, indexation, ASPA) sont réparties analytiquement
selon la structure de consommation par quintile ou la distribution des pensions. Pour les impôts sur les hauts revenus
et les patrimoines (M13), le rendement statique est corrigé d'une perte comportementale :
\begin{equation}
\text{perte} = e\,a\,\frac{\tau}{1-\tau}\times \text{rendement statique},\qquad e = 0{,}2~(0{,}1\text{--}0{,}4),\ a = 2
\end{equation}
pour les tranches supérieures de l'IR ; 0,5 pour l'élasticité des dividendes au PFU ; 25~\% puis 35~\% d'érosion pour
l'impôt sur la fortune ; 10~\% pour les droits de succession.

\paragraph{E. Calibration et validation.} Impôt net simulé 92,4\md{} contre 91,7\md{} publiés par la DGFiP
\preuve{A}. Barème non revalorisé : +2,7\md{} contre environ 4\md{} chiffrés officiellement (écart d'environ 30~\%).
RSA +10~\% : droits +2,8\md, versements +1,9\md.

\paragraph{F. Limites.} 3 millions de foyers hors champ ; revenus du capital et réductions d'impôt calés et non testés ;
éligibles à la prime d'activité sous-estimés d'environ 25~\% ; \textbf{la microsimulation est statique et ne voit pas
les salaires} : une hausse de salaire ne modifie pas les revenus de l'échantillon, d'où le recours à des cas types
pour étudier les effets de seuil (chapitre~\ref{ch:experiences}). Le dernier décile est mal représenté (rendement
quasi nul d'une hausse de la tranche supérieure de l'IR, voir chapitre~\ref{ch:audit}).

\section{Fiscalité des entreprises et cotisations sociales}
\label{sec:fiscent}

\paragraph{D. Équations.} Recettes d'impôt sur les sociétés :
$IS_t = IS_0\,\dfrac{EBE^m_t}{EBE^m_B}\,\dfrac{\tau_t}{\tau_B}$, avec $IS_0 = 15~\%$ des impôts courants sur le revenu
(ordre de grandeur) et $EBE^m$ l'excédent des branches marchandes. Cotisations : $D61_t = D61_B\,\dfrac{W_tN_t}{W_BN_B}\,\kappa_t$
où $\kappa_t$ traduit les mesures (une baisse de 10\md{} de cotisations employeurs diminue le coût du travail des
branches marchandes et la recette). Les cotisations employeurs entrent dans le coût salarial $CN_b$ des équations
d'emploi et de prix. Impôts de production : recette et choc de marge $mk$ (50~\% répercutés sur les prix de valeur
ajoutée, hypothèse \preuve{H}).

\paragraph{F. Limites.} Pas d'effet de l'IS sur la localisation des profits et des activités ; pas de distinction
entre PME et grandes entreprises ; la part de l'IS (15~\%) est un ordre de grandeur à caler sur les comptes de l'État.

\section{Prestations sociales, santé et autres dépenses publiques}
\label{sec:depenses}

\paragraph{A. Objet.} Dépenses publiques : rémunérations des agents, achats (consommation publique), prestations
sociales, investissement, subventions et intérêts.

\paragraph{D. Équations.}
\begin{align}
D1^{pub}_t &= D1^{pub}_B\,\frac{W_t}{W_B}\,\frac{N^{nm}_t}{N^{nm}_B}\,(1+\text{mesure}) \\
P3^{autres}_t &= (P3_B - D1^{pub}_B)\;\frac{G_t}{G_B}\;\frac{PC_t}{PC_B} \\
D62_t &= D62_B\;\frac{PC_{t-1}}{PC_B}\;\frac{V_t}{V_B} + \text{mesures} \\
\text{Autres}_t &= \alpha\,PIBN_t,\qquad \alpha = (TE_B - P3_B - D62_B - D41_B)/PIBN_B
\end{align}
La consommation publique en volume $G$ croît de 1,3~\% par an dans la référence (rythme 2011-2019) ; les « autres
dépenses » (subventions, investissement public, autres transferts) restent une part fixe du PIB nominal.

\paragraph{C. Hypothèse importante.} \textbf{La santé, l'éducation, la défense, la recherche et le fonctionnement
administratif ne sont pas distingués} : une économie d'un milliard d'euros sur l'une ou l'autre a exactement le même
effet macroéconomique (même contenu en importations, même effet direct sur la valeur ajoutée non marchande). Le modèle
ne mesure ni l'état de santé, ni le niveau d'éducation, ni la sécurité, ni leurs effets de long terme sur la
productivité. Il n'y a pas de module d'assurance maladie : les dépenses de santé sont incluses dans la consommation
publique et les prestations.

\section{Retraites (M9)}

\paragraph{D. Équations.} Pension à la liquidation de la génération $g$ :
$\text{pension}_g = \nu \times \overline{W}^{25}_g$ (moyenne des 25 dernières années de salaire moyen, revalorisées sur
les prix), puis revalorisation annuelle sur les prix (hors tabac) de l'année précédente, éventuellement diminuée d'une
sous-indexation. Le nombre de retraités par âge découle du profil d'activité (avant l'AOD, la moitié des sorties
précoces sont des départs anticipés ; après l'AOD, 40~\% des actifs cumulent emploi et retraite), calé sur Eurostat en
2015 (facteur 1,06 pour les retraités vivant à l'étranger). Lien avec M8 : population active, emploi des seniors
supplémentaires (89~\%) et masse des pensions.

\paragraph{E. Validation.} Masse des pensions reproduite à $\pm$2,3~\% sur 2008-2022 ; 2024 surestimée de 4,4~\%.
Suspension de la réforme de 2023 : 1,0\md{} en 2027 (estimations publiques : environ 1,4\md).

\paragraph{F. Limites.} Pas de régimes, de durée d'assurance, de décote ni de réversion (42\md) ; pas de solde du système
de retraite au sens du COR.

\section{Finances publiques, dette et taux d'intérêt (M8, M11)}
\label{sec:finpub}

\paragraph{D. Équations.} Recettes :
\begin{equation}
TR_t = \underbrace{D2_B\,\frac{C^{n}_t}{C^{n}_B}}_{\text{TVA, impôts sur produits}} + \underbrace{D5^{mén}_t + IS_t}_{\text{impôts sur le revenu}} + \underbrace{D61_t}_{\text{cotisations}} + \underbrace{(TR_B - D2_B - D5_B - D61_B)\,\frac{PIBN_t}{PIBN_B}}_{\text{autres}} + \text{mesures}
\end{equation}
Dépenses $TE_t = D1^{pub}_t + P3^{autres}_t + D62_t + D41_t + \text{Autres}_t$. Solde $S_t = TR_t - TE_t$. Dette :
\begin{equation}
Dette_t = Dette_{t-1} - S_t,\qquad D41_t = i^{app}_t\,Dette_{t},\qquad i^{app}_t = i^{app}_{t-1} + \tfrac18\bigl(\max(r_t,0) - i^{app}_{t-1}\bigr)
\end{equation}
Le taux apparent rejoint le taux de marché $r$ (OAT à 10 ans) au rythme d'un huitième par an (renouvellement de la
dette). Module européen (M11) : taux de la BCE $i_t = a + 0{,}73\,i_{t-1} + 0{,}39\,\pi^{ZE}_t + 0{,}30\,\text{écart}^{ZE}_t$
\preuve{E} ; transmission au taux long français $\Delta r_t = \alpha + 0{,}38\,\Delta i_t$ \preuve{E} ; prime de risque
$+3$ points de base par point de dette supplémentaire \preuve{H}.

\begin{exemple}[Exemple : l'« effet boule de neige »]
Avec une dette de 116~\% du PIB, un taux apparent qui monte de 2,2~\% (2025) vers 4,9~\% et une croissance nominale de
l'ordre de 3~\%, la dette augmente d'elle-même de $(i-g)\times d \approx (4{,}9-3)\times1{,}16 \approx 2{,}2$ points de
PIB par an, avant même le déficit primaire. C'est pourquoi la référence voit la dette dériver jusqu'à 308~\% du PIB en
2047 : ce chiffre n'est pas une prévision, mais la conséquence arithmétique de l'hypothèse d'un taux de 4,89~\% maintenu
sans aucune mesure de correction pendant vingt ans.
\end{exemple}

\paragraph{F. Limites.} Pas d'ajustement stock-flux en projection ; pas de structure par maturité de la dette ; pas de
sous-secteurs (État, sécurité sociale, collectivités) ; les ratios de dette dépendent fortement du niveau des prix
(dénominateur), lui-même dépendant de paramètres mal identifiés (courbe de Phillips, indexation).

\section{Inégalités, pauvreté et niveau de vie (M10)}
\label{sec:inegalites}

\paragraph{D. Indicateurs.} Sur l'échantillon pondéré par le nombre de personnes : médiane du niveau de vie, taux de
pauvreté (part des personnes sous 60~\% de la médiane), indice de Gini, rapports D9/D1 et S80/S20. Pour isoler l'effet
d'une mesure, le taux de pauvreté après réforme est calculé \textbf{au seuil de pauvreté d'avant réforme} (seuil fixe),
ce qui évite qu'une baisse générale des revenus « réduise » mécaniquement la pauvreté relative.

\paragraph{F. Limites.} Le niveau de vie est approché (revenu fiscal $-$ IR $+$ prestations simulées, par unité de
consommation) : il omet les revenus non imposables et les prestations non simulées. \textbf{Les niveaux ne sont pas
comparables aux statistiques de l'Insee} (médiane 22\,179~€ contre 26\,740~€ ; pauvreté 17,3~\% contre 15,4~\%), mais
les \emph{variations} dues à une mesure sont exploitables (chapitre~\ref{ch:reference}). Les effets sont statiques :
première année, sans changement de comportement ni effet macroéconomique en retour sur les revenus.

\section{Incertitude (M12) et sensibilité}

Les coefficients estimés sont tirés dans leur loi normale (matrice de variance-covariance robuste), les paramètres
calibrés uniformément dans leur plage (élasticité de substitution 0,2-0,6 ; effet du chômage sur les salaires
$-$0,8 à $-$0,2~\% ; fermeture d'offre 0,1-0,25). Le simulateur restitue l'intervalle entre le 5\textsuperscript{e} et le
95\textsuperscript{e} centile. La contre-expertise v0.6 ajoute un criblage de Morris sur 46 paramètres et des
sensibilités univariées (chapitre~\ref{ch:audit}). Non couvert : erreurs de données, de spécification, d'hypothèses
exogènes.

\section{Le catalogue de mesures : comment chaque levier entre dans le modèle}
\label{sec:catalogue}

Le tableau~\ref{tab:catalogue} indique, pour chaque curseur utilisé dans ce rapport, le canal par lequel il agit. C'est
là que se logent la plupart des hypothèses propres aux politiques publiques.

\begin{small}
\begin{longtable}{@{}p{3.2cm}p{7.2cm}p{3.6cm}@{}}
\caption{Traduction des leviers en entrées du modèle (fichier \code{simulateur/executer.py}).}\label{tab:catalogue}\\
\toprule
Levier & Canal de transmission & Paramètres clés \\
\midrule\endfirsthead
\toprule Levier & Canal de transmission & Paramètres clés \\ \midrule\endhead
\bottomrule\endfoot
IR (taux, seuils) & Microsimulation OpenFisca ; variation d'impôt par décile → RDB et consommation (M14) ; M13 pour les tranches supérieures & \preuve{A} barème ; $e = 0{,}2$ \\
CSG activité / retraite & Assiettes reconstituées sur l'échantillon, calées sur les totaux nationaux & \preuve{A/C} \\
TVA (normal, réduits) & Recette, prix à la consommation (80~\% répercutés), répartition par quintile de consommation & 8,2\md{} par point (normal) ; 5\md{} (réduits) \\
Taxe carbone & Recette et prix ; pas d'effet sur les émissions & 200\,M€ par €/t \\
PFU, ISF, successions & Recette corrigée des réactions (M13), revenu du dernier décile & 600\,M€/pt ; assiette ISF 930\md{} \\
IS & Coût d'usage du capital → investissement ; recette ; profits distribués (M14) & $\varepsilon = -0{,}18$ \\
Impôts de production & Recette ; 50~\% répercutés sur les prix ; profits (M14) & \preuve{H} \\
Cotisations employeurs & Coût du travail des branches marchandes (emploi, prix) et recette & \preuve{E/H} \\
SMIC & Diffusion à la masse salariale (0,11) ; coût public 0,66\md/pt & \preuve{H} \\
Salaires (hausse générale) & Choc sur l'équation de salaires & \preuve{E} \\
Point d'indice & Rémunérations publiques → RDB & \preuve{A} \\
Emploi public & Effectif non marchand (rampe de 5 ans), 56\,k€ par agent & \preuve{A} \\
Santé, éducation, défense, fonctionnement & Consommation publique (indistincte) & \preuve{H} \\
Investissement public & Demande + capital public productif ($e = 0{,}08$) & \preuve{L} \\
RSA, prime d'activité, AF, aides au logement & Microsimulation ; recours constant ; M14 & \preuve{A} \\
ASPA, indexation des pensions & Analytique ; M9 & \preuve{A} \\
Assurance chômage & Dépense ($-$40\md{} de base) ; chômage structurel $-$0,1 pt pour $-$10~\% & \preuve{H} (0 à 0,3) \\
Âge légal & M9 : actifs, emploi des seniors, pensions & \preuve{E} \\
Durée du travail & 70~\% de la variation d'heures absorbée par la productivité horaire ; salaires mensuels maintenus ; embauches compensatrices dans la fonction publique & \preuve{H} \\
Réindustrialisation & Parts importées de huit produits industriels (chimie, pharmacie, métallurgie, produits métalliques, électronique, équipements électriques, machines, automobile) réduites en 5 ans, après 2 ans de construction ; investissement de capacité ; surcoût de 10~\% & \preuve{H} (objectif, pas moyen) \\
Droits de douane & Prix des importations hors UE (45~\%), recette, riposte & \preuve{H} \\
Prix de l'énergie & Prix d'importation des produits B, C19, D & environnement \\
\end{longtable}
\end{small}
